\section{Five-coordinate normal-surface component certificates}
\label{nc:section}

Weighted orbit counting gives information about individual connected
components without expanding them.  We use this classical capability to
answer a specific question that the previous aggregate topology interface
could not answer: does a supplied, possibly disconnected normal surface
contain a compressing disk?  The contribution here is an explicit
constant-dimensional weight construction and its implementation contract.
Polynomial component analysis itself is already supplied by the weighted
Agol--Hass--Thurston algorithm \cite{AHT2006}.

\subsection{The validated input and its orbit system}
\label{nc:input}

Let $\mathcal T$ be a finite tetrahedral triangulation with $t$ tetrahedra
of a compact, connected, orientable $3$-manifold $M$, where $\partial M$
is one torus.  Face identifications within a tetrahedron and identifications
among its vertices or edges are permitted.  In particular, an edge may
have equal endpoints.  The input checks reciprocal face identifications,
edge orientations and links, finite vertex links, ambient orientability,
connectedness, and boundary topology.  Ideal or singular vertices are
outside this interface.  These checks establish neither irreducibility
nor that $M$ is the exterior of a specified knot diagram.

A nonnegative integral vector $x$ supplies seven normal coordinates per
tetrahedron: four triangle counts $T_0,T_1,T_2,T_3$ and three
quadrilateral counts $Q_0,Q_1,Q_2$.  The matching equations hold, and
at most one quadrilateral count in each tetrahedron is positive.  Thus
$x$ specifies a compact properly embedded normal surface $S$, possibly
empty, disconnected, or nonorientable.  Every integer is encoded in
binary, with signed hexadecimal transport available for large values.

Write $w_e=|S\cap e|$ for a global edge $e$.  If a local edge joins
vertices $a,b$, and $q(a,b)$ is the quadrilateral type avoiding this
edge, then its intersection count is
\[
 T_a+T_b+Q_0+Q_1+Q_2-Q_{q(a,b)}.
\]
The matching and gluing checks make this independent of the local
representative.  Choose an ordering and an orientation of the global
edges, and concatenate the $w_e$ intersection points into
\[
 X=\{0,\ldots,N-1\},\qquad N=\sum_e w_e.
\]
The points on $e$ occupy a consecutive block $I_e$.  Orientations are
transported through face gluings; they cannot be recovered just by
comparing global vertex labels, since those labels may coincide.

Each geometric triangular face is used once: select one representative
of an interior face pair and retain every boundary face.  Each of its
three normal-arc types gives an interval isometry between endpoint
blocks.  For the corner $v$ of a face opposite $f$, the number of arcs
is $T_v+Q_{q(f,v)}$.  Its endpoints occupy equally sized intervals
nearest $v$ on the two incident edges.  Depending on the transported
orientations, the isometry preserves or reverses order.  There are at
most $12t$ nonempty pairing families.  We retain their geometric
multiplicity, even when different families define identical relations.

\begin{lemma}\label{nc:orbit-components}
The orbits of these interval pairings are in bijection with the
connected components of $S$.
\end{lemma}
\begin{proof}
The normal cells give $S$ a finite cell structure.  Its vertices are
$X$, and its edges are the geometric normal arcs just described.
The pairing relation is exactly connectivity in this $1$-skeleton.
Each normal polygon has connected boundary in the $1$-skeleton, so
attaching it cannot join different graph components.  Conversely,
every polygon is incident to a vertex, and a path through polygons
can be replaced, for the purpose of connecting vertices, by paths
along their boundaries.  Thus the graph and surface components agree.
The argument permits loops, parallel edges, and repeated vertices on
a polygon boundary.
\end{proof}

\subsection{Two fixed cycles on the boundary torus}
\label{nc:boundary-basis}

Let $K$ denote the induced triangulation of $\partial M$, interpreted
as a cell structure with all edge occurrences retained.  Over
$\mathbb F_2$, put
\[
 Z_1=\ker\partial_1,\qquad B_1=\operatorname{im}\partial_2.
\]
Since $K$ is a torus, $\dim(Z_1/B_1)=2$.  The following elementary
construction fixes two graph cycles $\gamma_0,\gamma_1$ representing
a basis of this quotient.  It depends on $\mathcal T$, not on $x$.

Order the boundary edges and represent a chain by a bit vector.
Insert the boundary of every boundary triangle into a binary linear
basis.  Each occurrence of an edge is added by exclusive-or; two
occurrences therefore cancel.  This initially spans $B_1$.  Next
choose a spanning tree of the boundary graph.  If $P_u$ is the tree
path from its root to $u$, the fundamental cycle associated with a
non-tree edge $e=uv$ is
\[
 \gamma_e=P_u+e+P_v.
\]
For a loop, this is just $e$.  Parallel edges remain distinct.
Insert these fundamental cycles into the same linear basis, retaining
the original cycle whenever the insertion increases its rank.  Exactly
two are retained.  Indeed, the fundamental cycles span $Z_1$, and
the previously inserted triangle boundaries span $B_1$.  One may
use a spanning forest in the implementation without changing this
argument; the validated boundary graph is connected.

For a component $C$ of $S$, define the nonnegative integer
\[
 I_i(C)=\sum_{e\text{ in }\gamma_i}|C\cap e|,
 \qquad i=0,1,
\]
where an edge is in $\gamma_i$ precisely when its bit coefficient
is one.  These are integer intersection counts whose parities give
the mod-two pairings.  The cycles need not be simple or disjoint.
Keeping actual graph cycles, rather than just reduced quotient
vectors, makes their intersection weights directly available.

This basis choice does not introduce a topological dependence on
vertex or edge labels.  A different valid basis acts on the pair
of parities by an invertible matrix over $\mathbb F_2$, so their
simultaneous vanishing is unchanged.  The integer counts themselves
are relative to the chosen cycles and need not be invariant under
relabeling.  The implementation uses deterministic choices so that
its checker reconstructs the same field from the same input.  No
orientation of the cycles is needed, and no simplicity test for
them is part of the construction.

\subsection{A cellular weight field with five coordinates}
\label{nc:weights}

For a component $C$, let $F(C)$ be its number of normal polygon
pieces and let $B(C)=|C\cap K^{(1)}|$ be its number of vertices on
the ambient boundary.  Thus $B(C)$ is not the number of boundary
circles.  We construct a piecewise constant function
\[
 \lambda:X\longrightarrow\mathbb Z^5
\]
whose orbit sum is
\begin{equation}\label{nc:profile-vector}
 \bigl(\chi(C),F(C),B(C),I_0(C),I_1(C)\bigr).
\end{equation}
Individual Euler weights may be negative; signed additive weights
are essential to this construction.

Initially assign $(1,0,0,0,0)$ to every point.  For every original
geometric normal-arc family, subtract $(1,0,0,0,0)$ on its chosen
endpoint interval.  Using the domain interval after the pairing
object orders its two endpoint intervals is harmless: either
endpoint belongs to the same surface component.  This operation
must precede identity deletion, duplicate removal, or any orbit
simplification.  A loop still contributes one surface edge, and
two geometric arcs still contribute two edges even if their
connectivity relations agree.

For each normal polygon choose one corner, and add
$(1,1,0,0,0)$ there.  Entire parallel families are handled as
intervals.  A triangle at local vertex $v$ may use any incident edge
directed from $v$ towards another vertex: its chosen corners occupy
the initial interval of length $T_v$.  A quadrilateral family of
type $q$ may use any local edge $a<b$ with $q(a,b)\ne q$.
Measured from $a$, its corners occupy
\[
 [T_a,T_a+Q_q).
\]
They lie between the triangle stacks at the two endpoints.  There
is no other quadrilateral stack because the quadrilateral condition
has already been checked.

More generally, a local run $[r,r+h)$ measured from $a$ towards
$b$ becomes the same run in the global edge block when those
orientations agree, and becomes
\[
 [w_e-r-h,w_e-r)
\]
when they disagree.  Add the block's global offset in either case.
This prescription applies even when several local edges represent
one global edge.  The selected local corner is counted once for
its polygon, irrespective of other corners mapping to the same
global vertex.

Finally, on each boundary edge block $I_e$, add
\[
 (0,0,1,\mathbf 1_{e\in\gamma_0},\mathbf 1_{e\in\gamma_1}).
\]
An endpoint-event sweep adds these constant intervals and coalesces
adjacent equal vectors.  It retains zero-valued runs when necessary
to partition the entire universe.

\begin{lemma}\label{nc:cellular-weights}
For every orbit $O_C$ corresponding to a component $C$, the sum
$\sum_{p\in O_C}\lambda(p)$ equals
\eqref{nc:profile-vector}.
\end{lemma}
\begin{proof}
Every surface vertex receives one initial Euler contribution.
Every geometric surface edge contributes minus one at exactly one
of its endpoints, and every normal polygon contributes plus one
at exactly one corner.  All these incident points belong to the
same component as the cell being counted.  The first coordinate
therefore sums to $V(C)-E(C)+F(C)=\chi(C)$.  Polygon contributions
also give the second coordinate.  The remaining coordinates are
the indicators of boundary edge points and of membership in the
two selected cycle supports.  Their sums have exactly the stated
interpretations.
\end{proof}

The global sums of the first two coordinates can be checked
against the independently available total Euler characteristic
and total number of normal pieces.  These checks are useful
invariants, but componentwise correctness rests on the cellular
allocation in the lemma, not merely on agreement of global totals.

The allocation also avoids fractional local Euler weights.  One
could distribute a polygon's contribution among all its corners
or divide vertex contributions among incident tetrahedra, but
such formulas require extra incidence denominators.  Assigning
each cell to one incident vertex keeps every weight integral.
The individual values are not intended as intrinsic curvature:
changing a chosen endpoint or polygon corner changes the field
while preserving its sum on every component.  Hence repeated
incidences and large edge degrees affect neither the proof nor
the fixed number of coordinates.

\subsection{The disk and compressing-disk predicates}
\label{nc:disk-predicate}

Weighted AHT produces a compressed multiset
\[
 \{(m_j,z_j)\}_j,
 \qquad m_j\geq1,\quad z_j\in\mathbb Z^5,
\]
where $m_j$ orbits have feature vector $z_j$.  Equal vectors may
come from different components with different embeddings.  No
identification of those embeddings is needed for the following
predicates.

\begin{theorem}\label{nc:disk-criterion}
Suppose a component has profile $(c,f,b,u,v)$.  It is a disk if
and only if $c=1$ and $b>0$.  It is a compressing disk for
$\partial M$ if and only if, additionally,
\[
 (u\bmod2,v\bmod2)\ne(0,0).
\]
Consequently, summing the multiplicities of precisely these
profile groups counts the disk components and compressing-disk
components of the supplied surface.
\end{theorem}
\begin{proof}
A nonempty boundary of a normal surface component contains
normal arcs and hence intersects the boundary $1$-skeleton.
Thus $b>0$ is equivalent to nonempty boundary.  For an orientable
connected compact surface of genus $g$ with $r>0$ boundary
components, Euler characteristic one gives $2g+r=1$, hence
$g=0,r=1$.  For a nonorientable surface with $k\geq1$ crosscaps,
the same equation gives $k+r=1$, impossible when $r>0$.
The disk characterization follows.  Requiring boundary prevents
a closed projective plane from being classified as a disk.

Now let $C$ be such a disk.  Its boundary is a simple closed
curve on the torus.  Define a mod-two cochain on boundary edges by
$\alpha(e)=|C\cap e|\bmod2$.  Every boundary triangle has an even
number of arc endpoints, counted with occurrences, so $\alpha$
is a cocycle.  Its class is Poincar\'e dual to
$[\partial C]\in H_1(\partial M;\mathbb F_2)$.
Since $\gamma_0,\gamma_1$ form a homology basis, this class is
nonzero exactly when the two evaluations are not both zero.
These evaluations are $u\bmod2$ and $v\bmod2$.

An essential simple closed curve on a torus represents a
primitive nonzero class in integral homology, whose reduction
modulo two is nonzero.  An inessential simple closed curve bounds
a disk on the torus and has zero class.  Therefore the parity
test is equivalent to essentiality of $\partial C$, precisely
the definition of a compressing disk.  No irreducibility
assumption on $M$ is required.
\end{proof}

For arbitrary components with several boundary circles, their
mod-two classes can cancel.  The theorem therefore applies the
essentiality interpretation only after the disk test.  In
particular, it does not classify every boundary circle using
the aggregate pair of parities.

The distinction from an aggregate test is substantive even for
disks alone.  Consider two disjoint parallel meridian disks in a
solid torus together with an additional boundary-parallel disk.
The union is disconnected, has Euler characteristic three, and
has zero total boundary class modulo two.  Those aggregate data
do not identify a compressing disk.  Component weights retain
the two nonzero individual boundary classes, although their sum
vanishes, and the theorem counts both meridian disks.  Grouping
the two equal profiles into one record of multiplicity two
preserves exactly the information this predicate needs.  This
example explains why weighted orbit data add a capability that
cannot be recovered just from total Euler characteristic and
total boundary homology.

\begin{corollary}\label{nc:three-coordinates}
The three coordinates $(\chi(C),I_0(C),I_1(C))$ suffice for the
compressing-disk predicate.
\end{corollary}
\begin{proof}
A nonzero intersection parity implies $B(C)>0$.  Hence
$\chi(C)=1$ together with that parity already implies the full
criterion of Theorem~\ref{nc:disk-criterion}.
\end{proof}

The implementation retains five coordinates because they also
count inessential disk components and provide normal-piece and
boundary-size diagnostics.  No orientation-double query is
needed for these disk predicates.

\subsection{Binary complexity and certificate verification}
\label{nc:complexity}

\begin{lemma}\label{nc:linear-runs}
The initial field has $O(t)$ constant runs, independently of the
normal coordinates' magnitudes.  Each initial weight has
$\ell^1$ norm $O(t)$.
\end{lemma}
\begin{proof}
There is one initial interval, at most $12t$ arc-family
intervals, at most $5t$ positive polygon-family intervals, and
at most $6t$ boundary-edge intervals.  Each contributes at most
two endpoints.  For nonempty $X$, their common refinement thus
has at most $46t+1$ runs; cancellation or coincident endpoints
can only reduce this number.  Each interval contributes a
fixed vector of bounded norm, giving the second bound.
\end{proof}

All interval endpoints require $O(\log(N+1))$ bits.  The boundary
basis uses $O(t)$-bit vectors and ordinary binary elimination,
so its construction is polynomial in $t$.  Lemma~\ref{nc:linear-runs}
puts the weighted query within the polynomial bounds of the
classical AHT extension \cite{AHT2006}.  In particular, full-carrier
weight transfer increases the number of constant runs by at
most four per transfer; the algorithm does not allocate an
object for every represented point or every component.

The empty vector is handled separately by the same interface:
there are no orbit points, no weight runs, and no component
groups.  Its disk counts are zero.  For a nonempty vector,
$\log(N+1)$ is bounded by a constant multiple of the largest
coordinate bit length plus $\log(t+1)$, so this complexity
parameter is controlled by the actual input encoding.  A large
geometric number of normal pieces is therefore not itself an
obstruction to a short profile computation.

For signed weights, each accumulated value is a sum over
disjoint original points.  Its absolute norm is therefore at
most $N$ times the largest initial norm, giving
$O(\log(N+1)+\log(t+1))$ bits.  Together with the polynomial
number of orbit operations and weight breakpoints, this gives
a polynomial-time query in the encoded triangulation and
normal vector.  This is an application of the established
weighted component theorem, not a new general complexity
theorem for normal surfaces or unknot recognition.

The certificate checker reconstructs the validated geometry,
the two boundary cycles, and the five-coordinate field from
the supplied source input.  It checks the input fingerprint,
validates the local relations of the unweighted orbit trace,
and replays the additive transfers before comparing the
claimed profile multiset.  The fingerprint provides an identity
check; the mathematical conclusion comes from reconstruction
and replay.  The producer's search scheduler is not rerun.
Geometry and weight construction are shared with the producer,
so this is not a claim of a wholly independent geometric
implementation.

Verification is polynomial in the encoded input and supplied
trace length.  Exhausting a producer's cycle allowance gives
an inconclusive result, with no component verdict.  The
allowance does not bound all replay work or elapsed time;
cancellation is a separate cooperative callback whose
exceptions must propagate through validation.

Finally, a feature profile does not supply a component's full
normal coordinates or an explicitly embedded disk.  Recovering
such an object requires further coordinate weights or another
reconstruction mechanism.  A negative profile answer only
says that this particular $S$ has no compressing-disk component.
The existence search over normal surfaces, the relation to a
knot diagram, compressed three-dimensional cutting, and a
short global hierarchy remain separate obligations.
